% Version: v7 (2026-10-17), the Introduction explains what zero product determined algebras are used for, with Bresar Theorem 7.2 recalled, Remark rem:preserving carries it over to a bilinear map, and the open question on M_2(F_3) is answered by computation (Proposition prop:matrix). Acknowledgments removed. Claims added: Corollary cor:triangular and Proposition prop:matrix.
\documentclass[]{interact}
\usepackage[utf8]{inputenc}

\usepackage{mathtools}
\usepackage{amsmath,amssymb,amsthm}
\usepackage{amsfonts}
\usepackage{enumerate}
\usepackage{enumitem}

\usepackage[hidelinks,pagebackref=false]{hyperref}
\usepackage[all]{hypcap}
\hypersetup{linktoc=all}

\usepackage{quiver}
\usepackage{tikz-cd}
\usepackage[yyyymmdd,hhmmss]{datetime}
\usepackage{float}

\usepackage[numbers,sort&compress]{natbib}
\bibpunct[, ]{[}{]}{,}{n}{,}{,}
\renewcommand\bibfont{\fontsize{10}{12}\selectfont}
\makeatletter
\def\NAT@def@citea{\def\@citea{\NAT@separator}}
\makeatother

\theoremstyle{plain}
\newtheorem{theorem}{Theorem}[section]
\newtheorem{lemma}[theorem]{Lemma}
\newtheorem{corollary}[theorem]{Corollary}
\newtheorem{proposition}[theorem]{Proposition}

\theoremstyle{definition}
\newtheorem{definition}[theorem]{Definition}
\newtheorem{example}[theorem]{Example}
\newtheorem{question}[theorem]{Question}

\theoremstyle{remark}
\newtheorem{remark}[theorem]{Remark}
\newtheorem{notation}{Notation}

\allowdisplaybreaks

\newcommand{\Der}{\mathrm{Der}}
\newcommand{\Inn}{\mathrm{Inn}}
\newcommand{\Hom}{\operatorname{Hom}}
\newcommand{\JDer}{\mathrm{JDer}}
\newcommand{\BSDer}{\mathrm{BSDer}}
\newcommand{\ADer}{\mathrm{Der}_\alpha}
\newcommand{\LDer}{\mathrm{LDer}}
\newcommand{\Zs}{\mathcal{Z}}

\newcommand{\paperversion}{2}


\begin{document}
\title{Algebras Determined by a Braided-Symmetric Product}

\author{
\name{Abderrahim Adrabi\textsuperscript{a,b}\thanks{CONTACT A. Adrabi. Email: abderrahim.adrabi@um5r.ac.ma}}
	\affil{\textsuperscript{a}Faculty of Sciences, Mohammed V University in Rabat, Morocco;\\ \textsuperscript{b}Faculty of Educational Sciences, Mohammed V University in Rabat, Morocco}
}

\history{Preprint version \paperversion, compiled \today}

\maketitle

\begin{abstract}
An algebra is zero product determined when every bilinear map that vanishes on the pairs with zero product factors through the product.
We replace the product by any bilinear map $\mu$ and ask the same question.
For the symmetrized products of braided-symmetric derivations, this asks when the algebra with that product is zero product determined.
The answer is a quotient criterion that needs no surjectivity hypothesis.
For the symmetrized products of three braidings on a commutative algebra, $\mu$ has the form $f(xy)$ with $f$ linear.
The same holds on any algebra when the twist reverses products.
Every unital zero product determined algebra is determined by every such $\mu$.
The converse holds when $f$ is injective, so the identity braiding gives nothing new in characteristic other than $2$.
The converse fails for the twisted braiding, because the dual numbers are determined by their twisted product but are not zero product determined.
A computation over $\mathbb{F}_3$ shows that the algebra of $2 \times 2$ matrices is determined by a twisted product that has no form $f(xy)$.
\end{abstract}

\begin{keywords}
zero product determined algebra; braided-symmetric product; Jordan product; dual numbers
\end{keywords}

\begin{amscode}
16W10; 15A86; 18M05
\end{amscode}

\section{Introduction}

This section explains the question, states the answers and describes the plan.
Let $K$ be a field and let $A$ be a unital associative $K$-algebra.
Many properties of an algebra are decided by the pairs of elements whose product is zero \citep[Preface]{Bresar2021}.
This idea led to the zero product determined (zpd) algebras \citep[Definition 1.1]{Bresar2021}.
The notion arose from the study of commutativity preserving maps and of local derivations \citep[Preface]{Bresar2021}.
In both studies, results on linear maps came from bilinear maps that vanish on the pairs with zero product \citep[Preface]{Bresar2021}.
A basic case is a linear map $T \colon A \to B$ that preserves zero products, which means that $xy = 0$ implies $T(x)T(y) = 0$.
Such maps have been studied since Banach and Lamperti \citep[\S7.1]{Bresar2021}.
If $A$ is zpd, then $T(x)T(y) = S(xy)$ for a linear map $S$ \citep[Theorem 7.2]{Bresar2021}.
If $A$ and $B$ are unital and the range of $T$ contains an invertible element, then $T$ is an invertible multiple of a homomorphism \citep[Theorem 7.2]{Bresar2021}.
The same idea shows that every local derivation of a unital zpd associative algebra is a derivation \citep[Example 1.10]{Bresar2021}.
Jordan homomorphisms appear in the study of maps that preserve zero products or commutativity \citep[\S7.3]{Bresar2021}.
The Jordan product $x \circ y = xy + yx$ has its own class, the zero Jordan product determined (zJpd) algebras \citep[Definition 3.14]{Bresar2021}.

The product and the Jordan product are two choices of a bilinear map.
Bre\v{s}ar defines zpd for nonassociative algebras \citep[Definition 1.1]{Bresar2021}, so any bilinear map $\mu$ on $A$ gives a nonassociative algebra $(A,\mu)$.
We say that $A$ is \emph{determined by} $\mu$ when $(A,\mu)$ is zpd (Definition~\ref{def:determined}).
The main purpose of this paper is to decide which algebras are determined by which $\mu$.
The maps $\mu$ that interest us are the symmetrized products of \citet[Section 1]{Adrabi2026}.
Let $m \colon A \otimes A \to A$ be the linear map with $m(x \otimes y) = xy$.
In this paper a \emph{braiding} is a linear map $\sigma$ of $A \otimes A$, and we do not require $\sigma^2 = \mathrm{id}$, which differs from \citep[Definition 3.1]{Adrabi2026}.
Its symmetrized product is $x \cdot_\sigma y = \mu_\sigma(x,y) = m(x \otimes y + \sigma(x \otimes y))$.
We use three braidings.
The first is the identity. % claim-ok: definition of the first braiding
The second is the swap $x \otimes y \mapsto y \otimes x$. % claim-ok: definition of the second braiding
The third is the \emph{twisted} swap $x \otimes y \mapsto \alpha(y) \otimes \alpha(x)$, for a linear map $\alpha \colon A \to A$, which twists the swap \citep[Section 4]{Adrabi2026}.
Their symmetrized products are $2xy$, $xy + yx$ and $xy + \alpha(y)\alpha(x)$. % claim-ok: the formulas follow from the definitions given here
For an involution $\sigma$, a \emph{braided-symmetric derivation} is a linear map $D$ with $D(x \cdot_\sigma y) = D(x) \cdot_\sigma y + x \cdot_\sigma D(y)$ \citep[Definition 3.1]{Adrabi2026}.
These maps are the derivations of the nonassociative algebra $(A, \cdot_\sigma)$ \citep[Remark 3.4]{Adrabi2026}.
So an algebra determined by $\mu_\sigma$ is one whose braided-symmetric derivations can be studied through zero products, as local derivations are \citep[Example 1.10]{Bresar2021}.
Remark~\ref{rem:derivations} takes the first step.
For $T$ with $\mu(x,y) = 0 \Rightarrow T(x)T(y) = 0$, Remark~\ref{rem:preserving} gives the first part of Theorem~\ref{thm:bresar-zpp} on an algebra determined by $\mu$.

We now describe the main results.
Our first result is a criterion, which is Proposition 1.3 of Bre\v{s}ar for the algebra $(A,\mu)$.
It says that $A$ is determined by $\mu$ if and only if the tensors $x \otimes y$ with $\mu(x,y) = 0$ span the kernel of the induced map on $A \otimes A$ (Corollary~\ref{cor:criterion}).
Our second result concerns products of the form $\mu(x,y) = f(xy)$ with $f$ linear.
Every unital zpd algebra is determined by every such $\mu$ (Theorem~\ref{thm:zpd-fxy}).
So the zpd algebras of Bre\v{s}ar, such as matrix algebras and triangular algebras \citep[Corollary 2.17]{Bresar2021}, are determined by all these products (Corollary~\ref{cor:triangular}).
The converse holds when $f$ is injective (Proposition~\ref{prop:injective}).
The identity braiding has this form, and so does the twisted braiding on a commutative algebra (Corollary~\ref{cor:twisted}) or when $\alpha$ reverses products (Corollary~\ref{cor:anti}). % claim-ok: follows from the definitions just given
Our third result shows that the converse fails for non-injective $f$.
Theorem~\ref{thm:codim} gives a sufficient condition, and Example~\ref{ex:dual} shows that the dual numbers are determined by the twisted product but are not zpd.
So the twisted product describes maps on algebras that the zpd theory does not reach.
Our fourth result treats a twisted product that has no form $f(xy)$.
Proposition~\ref{prop:matrix} shows by computation that $M_2(\mathbb{F}_3)$ is determined by $xy + \alpha(yx)$, for conjugation $\alpha$ by $\operatorname{diag}(1,-1)$.
Question~\ref{q:noncomm} asks for the general case.

We follow this order.
Section~\ref{sec:prelim} recalls the definitions and the results of Bre\v{s}ar that we use, and it contains no new results.
Section~\ref{sec:criterion} proves the criterion.
Section~\ref{sec:fxy} treats products of the form $f(xy)$, with matrix and triangular algebras as applications.
Section~\ref{sec:beyond} contains the dual numbers, the matrix computation and the open question.

\section{Preliminaries}\label{sec:prelim}

This section recalls the notions of Bre\v{s}ar that we need, and none of it is new.
We recall two definitions and then the equivalent forms of the first.

\begin{definition}[{\citep[Definition 1.1]{Bresar2021}}]\label{def:bresar-zpd}
A nonassociative algebra $A$ over a field $F$ is said to be \emph{zero product determined} if, for every bilinear functional $\varphi \colon A \times A \to F$ satisfying $\varphi(x,y) = 0$ whenever $x, y \in A$ are such that $xy = 0$, there exists a linear functional $\tau$ on $A$ such that $\varphi(x,y) = \tau(xy)$ for all $x, y \in A$.
\end{definition}

We abbreviate \emph{zero product determined} to zpd, as \citet[Definition 1.1]{Bresar2021} does.

\begin{definition}[{\citep[Definition 3.14]{Bresar2021}}]\label{def:bresar-zjpd}
An associative algebra $A$ is \emph{zero Jordan product determined} if the Jordan algebra $A^+$ is zero product determined.
\end{definition}

Here $A^+$ has the Jordan product $x \circ y = xy + yx$ in place of $xy$. % claim-ok: the Jordan algebra A^+ of the cited definition

The next result gives equivalent forms of the first definition.

\begin{proposition}[{\citep[Proposition 1.3]{Bresar2021}}]\label{prop:bresar-zpd}
Let $A$ be a nonassociative algebra over a field $F$.
The following conditions are equivalent.
(i) $A$ is zpd.
(ii) If a bilinear functional $\varphi \colon A \times A \to F$ satisfies $\varphi(x,y) = 0$ whenever $xy = 0$, then $\sum_i \varphi(x_i, y_i) = 0$ whenever $\sum_i x_i y_i = 0$.
(iii) If $\sum_i z_i w_i = 0$, then $\sum_i z_i \otimes w_i$ lies in $\operatorname{span}\{x \otimes y \in A \otimes A \mid xy = 0\}$.
(iv) For every vector space $X$ over $F$ and every bilinear map $\varphi \colon A \times A \to X$ with $\varphi(x,y) = 0$ whenever $xy = 0$, there is a linear map $T \colon A \to X$ with $\varphi(x,y) = T(xy)$ for all $x, y \in A$.
\end{proposition}

The next definition recalls the derivations of the symmetrized product $x \cdot_\sigma y = \mu(x \otimes y) + \mu(\sigma(x \otimes y))$, where $\mu$ is the product of $A$. % claim-ok: linking sentence

\begin{definition}[{\citep[Definition 3.1]{Adrabi2026}}]\label{def:adrabi-bsder}
Let $A$ be a $K$-algebra and let $\sigma$ be a $K$-linear involution of $A \otimes A$.
A \emph{braided-symmetric derivation} for $\sigma$ is a $K$-linear map $D \colon A \to A$ that satisfies $D(x \cdot_\sigma y) = D(x) \cdot_\sigma y + x \cdot_\sigma D(y)$ for all $x, y \in A$.
\end{definition}

The next definition and theorem show what the zpd property is used for. % claim-ok: linking sentence
In them, algebras are associative \citep[\S7.1]{Bresar2021}.

\begin{definition}[{\citep[Definition 7.1]{Bresar2021}}]\label{def:bresar-zpp}
A linear map $T$ from an algebra $A$ to an algebra $B$ is said to be \emph{zero product preserving} if, for all $x, y \in A$, $xy = 0 \implies T(x)T(y) = 0$.
\end{definition}

The next theorem uses this notion. % claim-ok: linking sentence

\begin{theorem}[{\citep[Theorem 7.2]{Bresar2021}}]\label{thm:bresar-zpp}
Let $A$ and $B$ be algebras, and let $T \colon A \to B$ be a zero product preserving linear map. If $A$ is zpd, then there exists a linear map $S \colon A \to B$ such that $S(xy) = T(x)T(y)$ $(x, y \in A)$.
Moreover, if both $A$ and $B$ are unital and the range of $T$ contains an invertible element of $B$ (in particular, if $T$ is surjective), then $c = T(1)$ is invertible in $B$ and there exists a homomorphism $\phi \colon A \to B$ such that $T(x) = c\phi(x) = \phi(x)c$ $(x \in A)$.
\end{theorem}

\section{Determined by a bilinear map}\label{sec:criterion}

This section gives the definition and the criterion that later proofs use.
Throughout, $A$ is a unital associative algebra over a field $K$, and $\mu \colon A \times A \to A$ is $K$-bilinear.

\begin{definition}\label{def:determined}
The algebra $A$ is \emph{determined by} $\mu$ if the nonassociative algebra $(A,\mu)$ is zpd.
By \citet[Proposition 1.3(iv)]{Bresar2021} this holds if and only if, for every vector space $V$ and every bilinear map $\varphi \colon A \times A \to V$ with $\varphi(x,y) = 0$ whenever $\mu(x,y) = 0$, there is a linear map $\psi \colon A \to V$ with $\varphi(x,y) = \psi(\mu(x,y))$ for all $x, y \in A$.
\end{definition}

For $\mu(x,y) = xy$ this is the zpd property of $A$, and for $\mu = \circ$ it is the zJpd property of Definition~\ref{def:bresar-zjpd}.

The next remark carries the use of Theorem~\ref{thm:bresar-zpp} over to $\mu$.

\begin{remark}\label{rem:preserving}
Let $B$ be an algebra and let $T \colon A \to B$ be a linear map with $\mu(x,y) = 0 \Rightarrow T(x)T(y) = 0$.
If $A$ is determined by $\mu$, apply Definition~\ref{def:determined} to the bilinear map $(x,y) \mapsto T(x)T(y)$ with values in $B$.
It gives a linear map $S \colon A \to B$ with $S(\mu(x,y)) = T(x)T(y)$.
For $\mu(x,y) = xy$ this is the first part of Theorem~\ref{thm:bresar-zpp}.
\end{remark}

The next remark does the same for derivations. % claim-ok: linking sentence

\begin{remark}\label{rem:derivations}
Let $\sigma$ be an involution of $A \otimes A$ and let $D \colon A \to A$ be linear with $D(x) \cdot_\sigma y + x \cdot_\sigma D(y) = 0$ whenever $x \cdot_\sigma y = 0$.
If $A$ is determined by $\mu_\sigma$, apply Definition~\ref{def:determined} to the bilinear map $(x,y) \mapsto D(x) \cdot_\sigma y + x \cdot_\sigma D(y)$.
It gives a linear map $S \colon A \to A$ with $D(x) \cdot_\sigma y + x \cdot_\sigma D(y) = S(x \cdot_\sigma y)$.
So $D$ is a braided-symmetric derivation in the sense of Definition~\ref{def:adrabi-bsder} if and only if $S = D$ on the span of the products $x \cdot_\sigma y$.
\end{remark}


The criterion uses the linear map $\bar\mu \colon A \otimes A \to A$ with $\bar\mu(x \otimes y) = \mu(x,y)$.
Let $\Zs_\mu$ be the subspace of $A \otimes A$ spanned by the tensors $x \otimes y$ with $\mu(x,y) = 0$.
Then $\Zs_\mu \subseteq \ker \bar\mu$.

\begin{corollary}\label{cor:criterion}
The algebra $A$ is determined by $\mu$ if and only if $\Zs_\mu = \ker \bar\mu$.
\end{corollary}

\begin{proof}
Apply the equivalence of (i) and (iii) in Proposition~\ref{prop:bresar-zpd} to the nonassociative algebra $(A,\mu)$.
Condition (iii) says that $\sum_i z_i \otimes w_i \in \Zs_\mu$ whenever $\sum_i \mu(z_i,w_i) = 0$, which is $\ker \bar\mu \subseteq \Zs_\mu$.
\end{proof}

\section{Products of the form \texorpdfstring{$f(xy)$}{f(xy)}}\label{sec:fxy}

This section treats the case $\mu(x,y) = f(xy)$ with $f \colon A \to A$ linear, which includes two of the three braidings.

\begin{theorem}\label{thm:zpd-fxy}
Let $A$ be zpd, let $f \colon A \to A$ be linear, and let $\mu(x,y) = f(xy)$.
Then $A$ is determined by $\mu$.
\end{theorem}

\begin{proof}
Let $\varphi$ vanish on the pairs with $\mu(x,y) = 0$.
If $xy = 0$ then $\mu(x,y) = 0$, so $\varphi$ vanishes on zero products.
Since $A$ is zpd, there is a linear $\psi_0$ with $\varphi(x,y) = \psi_0(xy)$ by Proposition~\ref{prop:bresar-zpd}(iv).
Let $b \in \ker f$.
Then $\mu(b,1) = f(b) = 0$, so $\psi_0(b) = \varphi(b,1) = 0$.
Hence $\ker f \subseteq \ker \psi_0$, and $\psi_0$ induces a linear map on $A/\ker f$, which is isomorphic to the image of $f$.
This gives a linear map on the image of $f$, and we extend it to a linear map $\psi$ on $A$.
Then $\psi(f(z)) = \psi_0(z)$ for every $z$, so $\varphi(x,y) = \psi(\mu(x,y))$.
\end{proof}

The next proposition gives the converse for injective $f$.

\begin{proposition}\label{prop:injective}
Let $f \colon A \to A$ be linear and injective, and let $\mu(x,y) = f(xy)$.
If $A$ is determined by $\mu$, then $A$ is zpd.
\end{proposition}

\begin{proof}
Since $f$ is injective, $\mu(x,y) = 0$ if and only if $xy = 0$.
Let $\varphi$ vanish on zero products.
Then $\varphi$ vanishes on the pairs with $\mu(x,y) = 0$, and we get $\psi$ with $\varphi(x,y) = \psi(f(xy))$.
So $\varphi(x,y) = (\psi \circ f)(xy)$.
\end{proof}

The identity braiding gives $\mu(x,y) = xy + xy = 2xy$, so we take $f = 2\,\mathrm{id}$.

\begin{corollary}\label{cor:identity}
Let $\mu(x,y) = 2xy$.
\begin{enumerate}[label=(\roman*)]
\item If $A$ is zpd, then $A$ is determined by $\mu$.
\item If $\operatorname{char} K \neq 2$, then $A$ is determined by $\mu$ if and only if $A$ is zpd.
\end{enumerate}
\end{corollary}

\begin{proof}
Part (i) is Theorem~\ref{thm:zpd-fxy} with $f = 2\,\mathrm{id}$.
If $\operatorname{char} K \neq 2$ this $f$ is injective, and Proposition~\ref{prop:injective} gives the other direction of (ii).
\end{proof}

In characteristic $2$ the map $\mu$ is zero, every $\varphi$ in Definition~\ref{def:determined} is zero, and every algebra is determined by $\mu$.

The standard braiding gives the Jordan product, and determination by it is the zJpd property, as noted after Definition~\ref{def:determined}.
When $\operatorname{char} K \neq 2$ the map $q_\circ$ is onto, because $x \circ \tfrac12 = x$.
So Corollary~\ref{cor:criterion} recovers the usual form of the zJpd criterion.
It is not of the form $f(xy)$ on a noncommutative algebra, and we refer to \citet[Section 3.3]{Bresar2021} for its theory.

The twisted braiding has the form $f(xy)$ when $A$ is commutative. % claim-ok: shown by the next lines
Let $\alpha$ be an algebra endomorphism of $A$.
Then $\alpha(y)\alpha(x) = \alpha(yx) = \alpha(xy)$, so $xy + \alpha(y)\alpha(x) = (\mathrm{id} + \alpha)(xy)$.

\begin{corollary}\label{cor:twisted}
Let $A$ be commutative and zpd, and let $\alpha$ be an algebra endomorphism of $A$.
Then $A$ is determined by $\mu_\alpha(x,y) = xy + \alpha(y)\alpha(x)$.
\end{corollary}

\begin{proof}
Apply Theorem~\ref{thm:zpd-fxy} with $f = \mathrm{id} + \alpha$.
\end{proof}

The same holds on a noncommutative algebra when $\alpha$ reverses products. % claim-ok: shown by the next lines

\begin{corollary}\label{cor:anti}
Let $A$ be zpd, and let $\alpha \colon A \to A$ be linear with $\alpha(xy) = \alpha(y)\alpha(x)$ for all $x, y$.
Then $A$ is determined by $\mu_\alpha(x,y) = xy + \alpha(y)\alpha(x)$.
\end{corollary}

\begin{proof}
Here $\alpha(y)\alpha(x) = \alpha(xy)$, so $\mu_\alpha(x,y) = (\mathrm{id} + \alpha)(xy)$.
Apply Theorem~\ref{thm:zpd-fxy} with $f = \mathrm{id} + \alpha$.
\end{proof}

We use a result of Bre\v{s}ar on matrix algebras.

For a unital algebra $A_1$, a unital algebra $A_2$ and a unital $(A_1, A_2)$-bimodule $M$, the matrices $\left[\begin{smallmatrix} r_1 & m \\ 0 & r_2 \end{smallmatrix}\right]$ with $r_1 \in A_1$, $m \in M$ and $r_2 \in A_2$ form an algebra under the usual matrix operations, which Bre\v{s}ar denotes by $\mathrm{Tri}(A_1, M, A_2)$ \citep[\S2.1]{Bresar2021} \citep[Section 2.1]{Bresar2021}.

\begin{corollary}[{\citep[Corollary 2.17]{Bresar2021}}]\label{cor:bresar-matrix}
The following algebras are zpd:
(a) The algebra $M_n(B)$, where $n \geq 2$ and $B$ is any unital algebra.
(b) The algebra $\mathrm{Tri}(A_1, M, A_2)$, where $A_1$ and $A_2$ are unital algebras generated by idempotents and $M$ is a unital $(A_1, A_2)$-bimodule.
\end{corollary}

By Corollary~\ref{cor:bresar-matrix}(a) with $B = K$, the matrix algebra $M_n(K)$ is zpd for $n \geq 2$.
The transpose reverses products \citep[0.2.5]{HornJohnson2013}.
So Corollary~\ref{cor:anti} applies to $M_n(K)$ and to the transpose, and to every linear map $\alpha$ of $M_n(K)$ that reverses products.

Part (b) of the same corollary gives a second family, the triangular algebras.

\begin{corollary}\label{cor:triangular}
Let $A = \mathrm{Tri}(A_1, M, A_2)$ be as in Corollary~\ref{cor:bresar-matrix}(b), and let $f \colon A \to A$ be linear.
Then $A$ is determined by $\mu(x,y) = f(xy)$.
In particular $A$ is determined by $xy + \alpha(y)\alpha(x)$ for every linear $\alpha \colon A \to A$ with $\alpha(xy) = \alpha(y)\alpha(x)$.
\end{corollary}

\begin{proof}
By Corollary~\ref{cor:bresar-matrix}(b), $A$ is zpd, and $A$ has the unit with diagonal entries the units of $A_1$ and $A_2$.
Apply Theorem~\ref{thm:zpd-fxy}, and then Corollary~\ref{cor:anti} for the last claim.
\end{proof}

The algebra of upper triangular $n \times n$ matrices over a field is of type (b) \citep[Example 2.19]{Bresar2021}, so Corollary~\ref{cor:triangular} applies to it. % reproof-ok: application of the cited example

When $\alpha$ is an involution and $\operatorname{char} K \neq 2$, the map $\mathrm{id} + \alpha$ is twice the projection onto the fixed points of $\alpha$.
It is injective only for $\alpha = \mathrm{id}$, so Proposition~\ref{prop:injective} gives no converse.
The next section shows that the converse is false.

\section{Determined but not zpd}\label{sec:beyond}

This section gives a sufficient condition that does not assume zpd, and applies it to the dual numbers.

\begin{theorem}\label{thm:codim}
Let $f \colon A \to A$ be linear, let $N = \ker f$, and let $\mu(x,y) = f(xy)$.
Suppose that $f(1) \neq 0$, that $bn, nb \in N$ for all $n \in N$ and $b \in A$, and that there is a linear functional $c$ on $A$ with $x - c(x)1 \in N$ for every $x$.
Then $A$ is determined by $\mu$.
\end{theorem}

\begin{proof}
Let $\varphi$ vanish on the pairs with $\mu(x,y) = 0$.
If $n \in N$ then $ny$ and $xn$ lie in $N$, so $\mu(n,y) = \mu(x,n) = 0$ and $\varphi(n,y) = \varphi(x,n) = 0$.
Write $x = (x - c(x)1) + c(x)1$ and do the same for $y$.
Bilinearity gives $\varphi(x,y) = c(x)c(y)\varphi(1,1)$.
The same splitting gives $f(y) = c(y) f(1)$, and then $f(xy) = c(x) f(y) = c(x)c(y) f(1)$.
Since $f(1) \neq 0$, the map $a \mapsto a f(1)$ from $K$ to $A$ has a linear left inverse $g$.
Put $\psi(z) = g(z)\varphi(1,1)$.
Then $\psi(f(xy)) = c(x)c(y)\varphi(1,1) = \varphi(x,y)$.
\end{proof}

In words, the hypotheses say that $N$ is a two-sided ideal of codimension $1$ that does not contain $1$. % claim-ok: restatement of the hypotheses of the theorem
The next example meets them. % claim-ok: checked in the example

We also use the characterization of finite-dimensional zpd algebras by \citet[Theorem 2.27]{Bresar2021}.

\begin{theorem}[{\citep[Theorem 2.27]{Bresar2021}}]\label{thm:bresar-idempotents}
A finite-dimensional unital algebra is zpd if and only if it is generated by idempotents.
\end{theorem}

The next example uses this characterization to show that the class of algebras determined by the twisted product is strictly larger. % claim-ok: linking sentence for the example that follows

\begin{example}[Dual numbers]\label{ex:dual}
Let $\operatorname{char} K \neq 2$ and let $A = K[\varepsilon]/(\varepsilon^2)$, an algebra of the form studied in \citep[Example 2.23]{Bresar2021}, where a unity is adjoined to an algebra with trivial multiplication; here that algebra is $K\varepsilon$. % claim-ok: identification of the dual numbers with the algebra obtained from the one-dimensional algebra $K\varepsilon$ of trivial multiplication
The map $\alpha(a + b\varepsilon) = a - b\varepsilon$ is an involutive automorphism.
Since $A$ is commutative, $\mu_\alpha(x,y) = f(xy)$ with $f = \mathrm{id} + \alpha$, and $f(a + b\varepsilon) = 2a$.
So $N = K\varepsilon$, which is an ideal, $f(1) = 2 \neq 0$, and $c(a + b\varepsilon) = a$ satisfies the last hypothesis.
By Theorem~\ref{thm:codim}, $A$ is determined by $\mu_\alpha$.
The only idempotents of $A$ are $0$ and $1$, so $A$ is not generated by idempotents.
By Theorem~\ref{thm:bresar-idempotents}, $A$ is not zpd.
\end{example}

So the twisted braiding determines a class of algebras strictly larger than the zpd ones, even among commutative algebras.
When $\alpha$ is an automorphism of a noncommutative algebra, the twisted product $xy + \alpha(yx)$ is not of the form $f(xy)$, and Theorems~\ref{thm:zpd-fxy} and~\ref{thm:codim} do not apply.
The criterion of Corollary~\ref{cor:criterion} still applies.
The next proposition uses it on the smallest case. % claim-ok: linking sentence

\begin{proposition}\label{prop:matrix}
Let $A = M_2(\mathbb{F}_3)$ and let $\alpha$ be conjugation by $\operatorname{diag}(1,-1)$, an involutive automorphism.
Then $A$ is determined by $\mu_\alpha(x,y) = xy + \alpha(yx)$.
\end{proposition}

\begin{proof}
Here $A \otimes A$ has dimension $16$ over $\mathbb{F}_3$.
By Corollary~\ref{cor:criterion}, we must show $\Zs_{\mu_\alpha} = \ker \bar\mu_\alpha$, and the inclusion $\Zs_{\mu_\alpha} \subseteq \ker \bar\mu_\alpha$ always holds.
The matrix of $\bar\mu_\alpha$ in the bases of matrix units has rank $4$ over $\mathbb{F}_3$, so $\dim \ker \bar\mu_\alpha = 12$.
Among the $81 \cdot 81$ pairs $(x,y)$, exactly $417$ satisfy $\mu_\alpha(x,y) = 0$.
Gaussian elimination over $\mathbb{F}_3$ shows that the $417$ tensors $x \otimes y$ span a subspace of dimension $12$.
So the two subspaces have the same dimension and are equal.
\end{proof}

The same computation gives $\dim \Zs = \dim \ker = 12$ for the ordinary product, as expected since $M_2(\mathbb{F}_3)$ is zpd by Corollary~\ref{cor:bresar-matrix}(a), and $13 = 13$ for the twist by the transpose. % claim-ok: result of the same computation

\begin{question}\label{q:noncomm}
Which noncommutative algebras with an involutive automorphism $\alpha$ are determined by their twisted product $xy + \alpha(yx)$? % claim-ok: open question
\end{question}

\begin{thebibliography}{9}

\bibitem[Adrabi(2026)]{Adrabi2026}
A.~Adrabi, Braided-symmetric derivations of algebras, preprint, Mohammed V University in Rabat, 2026.

\bibitem[Bre\v{s}ar(2021)]{Bresar2021}
Bre\v{s}ar, M.: Zero Product Determined Algebras. Frontiers in Mathematics. Birkh\"auser, Cham (2021).

\bibitem[Horn and Johnson(2013)]{HornJohnson2013}
Horn, R.A., Johnson, C.R.: Matrix Analysis, 2nd edn. Cambridge University Press, Cambridge (2013).

\end{thebibliography}

\end{document}
